\section{Limitations}
\label{sec:limitations}

\textbf{The spatial model's geometry is measured; its corridor weights are not.} The
network, the yards and the route-kilometre audit are data. The service intensity assigned
to each corridor is declared judgment, because Pakistan Railways publishes no
service-frequency statistics and its public timetable interface refuses unauthenticated
access; the basis of every weight is stated on its data row, the freight-wagons proxy is
retained as the structural alternative, and the split's real uncertainty is the distance
between the two published alternatives. A timetable extraction, or divisional train-km if
any edition ever carries them, would replace the judgment with measurement. The national
total is unaffected throughout.

\textbf{Workshop and stationary fuel is inside the module's metered total, measured, and
booked at its own category.} Fuel burnt by railway workshops, power houses, pumps and track
machines is part of the same consumption total this sector reads---the source says so
itself: ``these quantities comprise the entire fuel consumed on the Railways, i.e.\ on
locomotives, power houses, water pumps, workshops etc.''---and IPCC §3.4.1.4 flags the
overlap explicitly. That fuel belongs outside 1A3c, but not in any single category. This
paper previously assigned the whole quantity to 1A4c~ii, which was a classification by
purchaser rather than by source: under the IPCC definitions, track-maintenance machinery is
off-road transport, workshop generators and power houses are stationary combustion, and
pumps follow their end use. What remains unresolved is only that FINAL split: the measured
non-traction quantity books at its own category (1A2/1A4c unresolved) rather than inside
1A3c, and no published record yet decomposes it across those end uses.

This paper previously also called that fuel inseparable. It is not. The same edition carries a
detailed fuel table on p.163 splitting locomotive from other-than-locomotive consumption, for
diesel and furnace oil, back to 1950---annually from 2015--16, in five-year averages before
that---one page from the summary this sector was built on. Non-traction diesel is
\SI{1599}{} thousand litres, \SI{1.21}{\percent} of the total (since the 2024--25
restatement below, the two tables' denominators agree). The
overlap is small, which is what the sector had assumed; the difference is that it is now
measured---and published as a ledger:
\texttt{outputs/railways\_fuel\_claimants\_2024.csv} carries each claimant with its
volume, share and emissions at the sector's own factors, closing to the headline exactly.
The recode IS the reference: the engine books
\SI{0.00498}{\mega\tonne} at its own non-traction category and 1.A.3.c at locomotive
traction, with the older summary booking preserved as the configuration
\emph{rail\_ie\_summary} in
\texttt{outputs/national\_configurations.csv}. Either booking conserves the national
total at \num{444.95542} exactly---the same fuel at the same
factors under a different code, never a deletion. (An earlier edition priced this axis
as a subtraction; the conserving recode replaced it, and the reference itself now
carries the split.)

\textbf{The Year Book disagreed with itself about the reference year, and its next edition
settled it.} Checking the summary against the detail across every year both tables cover
produced exactly one failure: the 2023--24 edition printed the FY2023--24 locomotive cell as
``136.988''---a decimal point where every other row in its column uses a comma---making the
detail sum \num{138587} against the summary's twice-stated \num{132588}, a
\SI{4.52}{\percent} gap. This paper reported the summary, on the strength of two statements
against one typographically anomalous row, and published the detail reading as a labelled
sensitivity of \SI{0.412}{\mega\tonne} against \SI{0.394}{}.

The 2024--25 edition resolved it. Its historical run restates the cell as
\num{130988}---printed with the same decimal-point quirk on the same row---and
$\num{130988}+\num{1599}=\num{132587}$, within one part in \num{132588} of the summary. The
misprint hypothesis was right, the reported number never moved, and the sensitivity is
retired. The reconciliation now closes in every year:

\begin{table}[htbp]
\centering
\footnotesize
\begin{tabular}{@{}lrrrr@{}}
\toprule
\textbf{FY} & \textbf{Summary} & \textbf{Loco} & \textbf{Other} & \textbf{Detail} \\
\midrule
2015--16 & 139\,487 & 137\,908 & 1\,579 & 139\,487 \\
2016--17 & 144\,631 & 142\,607 & 2\,024 & 144\,631 \\
2017--18 & 156\,396 & 154\,357 & 2\,039 & 156\,396 \\
2018--19 & 159\,927 & 157\,979 & 1\,948 & 159\,927 \\
2019--20 & 139\,600 & 137\,969 & 1\,630 & 139\,599 \\
2020--21 & 140\,257 & 138\,841 & 1\,416 & 140\,257 \\
2021--22 & 142\,086 & 140\,482 & 1\,601 & 142\,083 \\
2022--23 & 104\,909 & 103\,272 & 1\,637 & 104\,909 \\
2023--24 & 132\,588 & 130\,988 & 1\,599 & 132\,587 \\
2024--25 & 138\,673 & 137\,096 & 1\,577 & 138\,673 \\
\bottomrule
\end{tabular}
\caption{Thousand litres of diesel, across every year the two tables both cover. Eight
reconcile to the digit and two to within three parts in \num{140000}; the FY2023--24 row
carries the 2024--25 edition's restatement of the cell its predecessor misprinted as
``136.988.''}
\label{tab:yearbook}
\end{table}

The check that found this ran late, and the honest reflection stands: this paper's argument
is that a national submission can be refuted on a fuel account, and an inventory making that
argument has to hold its own fuel account to the same standard. It now does, on every
build---and the one anomaly the check ever surfaced was a misprint in a national statistical
publication, identified here from the document's own internal evidence before the publisher
restated it.

\textbf{The furnace-oil basis is settled, and it was not what the summary said.} Editions
changed the published unit label without announcement, and this paper recorded the ambiguity as
unresolved. The 2023--24 edition resolves it against itself: the p.163 column is headed
``Furnace Oil (M.~Tonnes)'', the narrative on the same page states \SI{79.907}{} M.T, and that
is exactly the detail table's \num{15.083} locomotive plus \num{64.824} other. The Principal
Statistics table calls the same quantity ``Liters/thousand''---and its unit column is
demonstrably unreliable, labelling Persons Employed ``Tonne'' and Cost of Employees ``No.'' The
basis is tonnes. The stream is under \SI{0.001}{\mega\tonne} either way, as this paper already
argued, so nothing material moves; the data now say the right thing.

\textbf{\NtwoO{} is a tenth of the sector and carries $\pm\SI{150}{\percent}$.} The rail
default's 95\% interval spans \SIrange{14.3}{85.8}{\kilo\gram\per\tera\joule}, which is the
whole of the sector's uncertainty in practice. Nothing short of measurement on Pakistani
locomotives would narrow it, and the sector is too small for that to be worth doing.

\textbf{The air-pollutant layer rests on European fleet-average factors and a declared
sulfur assumption.} The EMEP/EEA Tier~1 locomotive factors describe an average European
fleet, not Pakistan's; the chapter's own confidence intervals are carried on every factor
row, and its printed intervals for the particulate rows do not bracket their central
values---a defect of the source recorded on the rows rather than silently repaired. \SOtwo{}
depends entirely on fuel sulfur, for which no operator certificate is published: it is
reported as explicit sulfur cases spanning \SIrange{10}{10000}{ppm} rather than as one
number, with the pre-Euro-V \SI{500}{ppm} specification as the declared reference.

\textbf{Wear PM is exploratory and European-anchored.} The wheel, rail and brake-wear
layer (\citet{tno2024rail}'s fuel-based factors on the metered diesel: roughly
\SI{15}{\tonne} of PM\textsubscript{10} a year, a tenth of exhaust PM\textsubscript{10},
alongside it and never summed) rests on Dutch and Austrian measurements. Pakistani
conditions---tread brakes throughout, heavy axle loads, track condition---plausibly sit
higher, the source's own PM\textsubscript{2.5} bullet contradicts its stated size
fractions (recorded on the factor rows, the fraction-consistent reading adopted), and
$\pm$\SI{150}{\percent} is carried until a tread-brake-specific or Pakistani measurement
exists.

\textbf{One operator, one figure, no redundancy.} The estimate has no second source. The only
independent check the publication offers on itself is that the FY2023 figure is republished
identically in two consecutive editions, and the traffic-based intensity check of
Section~\ref{sec:validation}. If the published fuel figure were wrong, nothing here would
detect it.
